\section*{Chapter \ref{ch:b2:equilibrium-shifts} --- Equilibrium Constants and Shifts}

\begin{solution}{exo:b2:equilibrium-shifts:1}
$K^\circ = \exp(32\,800/(8.314 \times 298.15)) = \num{5.6e5}$, and for
$+\qty{32.8}{kJ/mol}$ its inverse, \num{1.8e-6}. A factor of 100 needs
$\Delta(\Delta_r G^\circ) = -RT\ln 100 = \qty{-11.4}{kJ/mol}$.
\end{solution}

\begin{solution}{exo:b2:equilibrium-shifts:2}
$\ln K^\circ(400) = \ln(5.4 \times 10^5) + (91\,880/8.314)(1/400 -
1/298.15) = 13.20 - 9.44 = 3.76$, $K^\circ \approx 43$. The tables give 35.6:
over \qty{100}{K} the constant-$\Delta_r H^\circ$ approximation is already
off by 20\,\% (the reaction becomes more \omterm{def:b2:reaction-enthalpy:exothermic}{exothermic} as $T$ rises).
\end{solution}

\begin{solution}{exo:b2:equilibrium-shifts:3}
(a) $n = 1$, $r = 0$, $\varphi = 2$: $v = 1$ (the vapour pressure is a function
of $T$). (b) $v = 3 - 1 + 2 - 3 = 1$. (c) $v = 3 - 1 + 2 - 1 = 3$. (d) Four
species, one reaction, two phases (carbon and the gas): $v = 4 - 1 + 2 - 2 =
3$; if the gas is made only from carbon and steam, $n(\ce{CO}) = n(\ce{H2})$
in the gas, $s = 1$ and $v = 2$.
\end{solution}

\begin{solution}{exo:b2:equilibrium-shifts:4}
Heating: \omterm{def:b2:reaction-enthalpy:exothermic}{endothermic} direction, towards \ce{NO2} (forward). Compressing:
towards fewer gas molecules, \ce{N2O4} (backward). Argon at constant volume:
no change, the partial pressures of the reacting gases are unchanged.
\end{solution}

\begin{solution}{exo:b2:equilibrium-shifts:5}
At \qty{298}{K}: $\Delta_r G^\circ = 180.6 - 298.15 \times 0.0248 =
\qty{173.2}{kJ/mol}$, $K^\circ = \num{4.5e-31}$. At \qty{2000}{K}: $\Delta_r
G^\circ = 180.6 - 2000 \times 0.0248 = \qty{131.0}{kJ/mol}$, $K^\circ =
\num{3.8e-4}$. With $\Delta\nu_{\text{gas}} = 0$, $x(\ce{NO})^2 = K^\circ x(\ce{N2})
x(\ce{O2})$: $x(\ce{NO}) = \sqrt{3.8 \times 10^{-4} \times 0.78 \times 0.21} =
\num{7.9e-3}$, almost one per cent. Engines burn at such temperatures; as the
gases cool, $K^\circ$ falls enormously, but the decomposition of \ce{NO} is
very slow: the hot equilibrium is frozen in the exhaust (until a catalyst
removes it).
\end{solution}

\begin{solution}{exo:b2:equilibrium-shifts:6}
With $\xi$ mol of ester and a constant total amount, $Q =
\xi^2/[(1 - \xi)(n_B - \xi)]$. For $n_B = 1.0$: $\xi^2/(1 - \xi)^2 = 4$, $\xi
= \qty{0.67}{mol}$. For $n_B = 3.0$: $3\xi^2 - 16\xi + 12 = 0$, $\xi =
\qty{0.90}{mol}$. Removing water keeps $Q$ below $K^\circ$ whatever $\xi$: the
reaction keeps going forward until the acid is consumed.
\end{solution}

\begin{solution}{exo:b2:equilibrium-shifts:7}
Without argon: amounts $1 - \alpha$, $2\alpha$, total $1 + \alpha$;
$Q = 4\alpha^2/(1 - \alpha^2) = 0.148$, $\alpha = \sqrt{0.148/4.148} = 0.189$.
With \qty{1}{mol} argon at \qty{1}{bar}: total $2 + \alpha$, $Q =
4\alpha^2/[(1 - \alpha)(2 + \alpha)] = 0.148$, so $4.148\alpha^2 + 0.148\alpha
- 0.296 = 0$, $\alpha = 0.250$: diluting at constant total pressure lowers the
partial pressures, which favours the side with more gas molecules. At
constant volume, the partial pressures of the reacting gases do not change:
$\alpha$ stays 0.189.
\end{solution}

\begin{solution}{exo:b2:equilibrium-shifts:8}
Only the solid: $n = 3$, $r = 1$, $s = 1$ ($p(\ce{NH3}) = p(\ce{HCl})$),
$\varphi = 2$: $v = 3 - 1 - 1 + 2 - 2 = 1$; choose $T$ and the pressure is
fixed (a ``dissociation pressure''). With added ammonia, $s = 0$ and $v = 2$:
$T$ and one pressure can be chosen; the product $p(\ce{NH3})\,p(\ce{HCl})$ is
fixed by $T$.
\end{solution}

\begin{solution}{exo:b2:equilibrium-shifts:9}
$\Delta_r H^\circ = -R\ln(K_2/K_1)/(1/T_2 - 1/T_1) = -8.314 \times
\ln(0.0173)/(1/600 - 1/500) = \qty{-101}{kJ/mol}$, more negative than at
\qty{298}{K} (\qty{-91.9}{kJ/mol}), as \omterm{thm:b2:reaction-enthalpy:kirchhoff}{Kirchhoff's law} predicted
($\Delta_r C_p^\circ < 0$).
\end{solution}

\begin{solution}{exo:b2:equilibrium-shifts:10}
$Q = \dfrac{n_{\ce{NH3}}^2\,n_{\text{tot}}^2}{n_{\ce{N2}}n_{\ce{H2}}^3}
\left(\dfrac{p^\circ}{p}\right)^2$. At constant $T$, $p$:
$\partial\ln Q/\partial n_{\ce{N2}} = -1/n_{\ce{N2}} + 2/n_{\text{tot}}$, positive
when $n_{\ce{N2}} > n_{\text{tot}}/2$, that is $x(\ce{N2}) > 1/2$. Then $Q$
exceeds $K^\circ$ and the equilibrium shifts backward. At constant volume,
$\partial\ln Q/\partial n_{\ce{N2}} = -1/n_{\ce{N2}} < 0$: always forward.
\end{solution}

\begin{solution}{exo:b2:equilibrium-shifts:11}
$K^\circ = \exp(-1620/(8.314 \times 1100)) = 0.84$: $p(\ce{CO2}) =
\qty{0.84}{bar}$ at equilibrium, that is $n = pV/RT = 0.84 \times 10^5 \times
0.0100/(8.314 \times 1100) = \qty{0.092}{mol}$ of gas. With \qty{0.050}{mol}
of limestone, all of it decomposes ($p = \qty{0.46}{bar} < K^\circ p^\circ$):
no equilibrium, a break of equilibrium. With \qty{0.200}{mol}, equilibrium:
\qty{0.092}{mol} of \ce{CO2} and \ce{CaO}, \qty{0.108}{mol} of \ce{CaCO3} left.
\end{solution}

\begin{solution}{exo:b2:equilibrium-shifts:12}
Bed 1 leaves at about \qty{890}{K} and 68\,\% conversion, bed 2 at about
\qty{785}{K} and 92\,\%, bed 3 at about \qty{725}{K} and 98\,\%. In one
adiabatic bed the gas heats up as it reacts until it meets the equilibrium
curve, which at that temperature allows only about 68\,\%. A cold catalyst
throughout is not an option because the catalyst works only when hot; cooling
between beds keeps the rate and moves the equilibrium limit up. The last bed
is limited by the equilibrium at the lowest temperature at which the catalyst
still works (plants then absorb the trioxide and pass the gas once more over
a catalyst).
\end{solution}

\begin{solution}{pb:b2:equilibrium-shifts:1}
\textbf{1.} $\Delta_r H^\circ = \qty{-91.88}{kJ/mol}$, $\Delta_r S^\circ =
2(192.77) - 191.61 - 3(130.68) = \qty{-198.1}{J/(K.mol)}$, $\Delta_r G^\circ =
\qty{-32.8}{kJ/mol}$, $K^\circ = \num{5.6e5}$.
\textbf{2.} $K^\circ(700) = \exp(-54\,380/(8.314 \times 700)) = \num{8.8e-5}$;
$K^\circ(800) = \exp(-77\,320/(8.314 \times 800)) = \num{8.9e-6}$.
\textbf{3.} $\Delta_r H^\circ = -R\ln(K_{800}/K_{700})/(1/800 - 1/700)$,
that is \qty{-106}{kJ/mol}, more negative than at \qty{298}{K}.
\textbf{4.} $\ln K^\circ(723) = \ln(8.75 \times 10^{-5}) + 12\,777 \times
(1/723.15 - 1/700)$, with $12\,777 = 106\,230/8.314$: $K^\circ =
\num{4.9e-5}$.
\textbf{5.} $\Delta_r G^\circ(723) = -91.88 + 723.15 \times 0.1981 =
\qty{51.4}{kJ/mol}$, $K^\circ = \num{1.9e-4}$: four times too large, because
$\Delta_r H^\circ$ and $\Delta_r S^\circ$ change over \qty{425}{K}
($\Delta_r C_p^\circ \approx \qty{-44}{J/(K.mol)}$).
\textbf{6.} \ce{N2} $1 - \xi$, \ce{H2} $3 - 3\xi$, \ce{NH3} $2\xi$, total $4 -
2\xi$.
\textbf{7.} $x = 2\xi/(4 - 2\xi)$, so $1 - x = (4 - 4\xi)/(4 - 2\xi)$; the
nitrogen fraction $(1 - \xi)/(4 - 2\xi) = (1 - x)/4$ and the hydrogen fraction
is three times larger.
\textbf{8.}
\[
  K^\circ = \frac{x^2}{\frac{1-x}{4}\left(\frac{3(1-x)}{4}\right)^3}
  \left(\frac{p^\circ}{p}\right)^2 = \frac{256}{27}\,\frac{x^2}{(1 - x)^4}
  \left(\frac{p^\circ}{p}\right)^2 ;
\]
take the square root.
\textbf{9.} $a = (\sqrt{27}/16) \times \sqrt{4.9 \times 10^{-5}} \times 1 =
\num{2.27e-3}$; $x/(1 - x)^2 = a$ gives $x = 0.0023$.
\textbf{10.} $a = 0.455$; $ax^2 - (2a + 1)x + a = 0$: $x = 0.25$.
\textbf{11.} $n = 3$, $r = 1$, $s = 1$ (ratio 1 : 3 kept by the reaction),
$\varphi = 1$: $v = 2$. Once $T$ and $p$ are chosen, nothing is left free.
\textbf{12.} Pressure: forward ($\Delta\nu_{\text{gas}} = -2$). Temperature:
backward (\omterm{def:b2:reaction-enthalpy:exothermic}{exothermic}).
\textbf{13.} An inert gas at constant total pressure lowers the partial
pressures of the reacting gases, as a pressure drop would: backward.
\textbf{14.} $\partial\ln Q/\partial n_{\ce{N2}} = (-1/0.30 + 2)/n_{\text{tot}} < 0$:
$Q$ falls below $K^\circ$, forward.
\textbf{15.} No: at constant volume $\partial\ln Q/\partial n_{\ce{N2}} =
-1/n_{\ce{N2}} < 0$ whatever the composition: forward.
\textbf{16.} With almost no ammonia, $Q$ is far below $K^\circ$ at the inlet:
the gas enters far from equilibrium and reacts forward.
\textbf{17.} The gas spends too short a time on the catalyst for the
equilibrium to be reached; a longer contact would cost a larger, more
expensive converter for a small gain.
\textbf{18.} $x = 2\xi/(4 - 2\xi) = 0.18$ gives $\xi = \qty{0.305}{mol}$ per
mole of \ce{N2} fed: \omterm{def:b2:equilibrium-shifts:single-pass}{single-pass conversion} of hydrogen (and of nitrogen)
$0.305$, 31\,\%.
\textbf{19.} In steady operation all the fresh feed is eventually converted:
\qty{2.00}{mol} of ammonia leave per unit time; the converter must pass $1.00/
0.305 = \qty{3.3}{mol}$ of \ce{N2} per unit time, the rest being recycled.
\textbf{20.} The argon (and methane) that enter with the feed never react and
are not removed with the liquid ammonia: without a \omterm{def:b2:equilibrium-shifts:single-pass}{purge} they would
accumulate in the loop and lower the partial pressures of the reactants.
\textbf{21.} At \qty{500}{K} the iron catalyst is too slow: the equilibrium
would be favourable but would never be approached in an industrial time.
\textbf{22.} At \qty{723}{K} and \qty{200}{bar}, stoichiometric feed, perfect
gases: $\boldsymbol{x(\ce{NH3}) \approx 0.25}$.
\end{solution}
